\section{What Is a Field? From Particles to Continuous Degrees of Freedom}

A particle model assigns a finite amount of data to each particle: position,
velocity, momentum, charge, spin, or other variables. A field model assigns data
to each point of a continuous domain. For example, a temperature field assigns a
number to each point in a body, while an electromagnetic field assigns electric
and magnetic vectors to each point of space and time.

The move from particles to fields is a move from finitely many degrees of
freedom to continuously many degrees of freedom. This is why field theory is
mathematically richer than point-particle mechanics. It involves functions,
partial differential equations, boundary conditions, and local conservation
laws.

The basic question is always:
\[
    \text{What physical quantity is assigned to each point?}
\]
For a scalar field, the answer is one number. For a vector field, it is a
vector. For a tensor field, it is a multilinear object whose components depend
on coordinates but whose meaning should not depend on the chosen coordinate
system.

\begin{definitionbox}
A \emph{classical field} is a physical quantity represented by an assignment of
mathematical data to each point of space, spacetime, or another continuous
domain. The value assigned at a point may be a scalar, vector, tensor, spinor,
or more general geometric object.
\end{definitionbox}

The word \emph{classical} means that the field is treated as a definite object,
not as an operator or quantum state. Classical field theory studies equations
for such fields before quantization.
